% =====================================================================
\section{Diode characterization: I--V and C--V measurements}\label{sec:diodes}
% =====================================================================
% PAGE BUDGET for the whole section: about 3 pages (figures included).
% Sections 2.1 to 2.5 are written by the I-V group. Section 2.6 (C-V) is written by us.
\subsection{Theoretical Introduction}

\placeholder{}
\subsection{Experimental setup and current resolution}\label{sec:iv-setup}


For this part of the practical we observe the behavior of a p-n junction using two diodes, one square and one rake-shaped, under illumination and in the dark.
Both junctions are set up on a silicon wafer. In this setup the wafer is positioned under a microscope. The SMU is connected through tips on the cathode side and
a ground plate from the microscope on the anode side.
We observed that the forward biased region lies in the negative voltages because the n-type material is on the top side. Because of this, we use a sweep from $\SI{-5}{\V}$ to $\SI{10}{\V}$ in steps of $\SI{100}{\milli \V}$ for the measurements. 
\begin{figure}[h!]
  \centering
  \begin{subfigure}[b]{0.48\linewidth}
    \centering
    \begin{tikzpicture}[>={Stealth[length=2mm]}, font=\small]

  % layers
  \fill[blue!15]   (0,1.25) rectangle (5,2.0);     % n-type
  \fill[yellow!40] (0,0.75) rectangle (5,1.25);    % depletion region
  \fill[red!15]    (0,0)    rectangle (5,0.75);    % p-type
  \draw (0,0) rectangle (5,2.0);
  \draw[dashed] (0,1.0) -- (5,1.0);                % junction

  % contacts
  \filldraw[fill=gray!40] (0,2.0)  rectangle (5,2.25);
  \filldraw[fill=gray!40] (0,-0.25) rectangle (5,0);

  % labels
  \node at (2.5,1.625) {n};
  \node at (2.5,0.375) {p};
  \node[font=\footnotesize] at (2.5,1.0) [fill=yellow!40,inner sep=1pt] {depletion};

  % terminals
  \draw (2.5,2.25) -- (2.5,2.9) node[above] {Cathode};
  \draw (2.5,-0.25) -- (2.5,-0.9) node[below] {Anode};

\end{tikzpicture}
    \caption{Schematic of the wafer cross-section\protect\footnotemark}
    \label{fig:junction}
  \end{subfigure}\hfill
  \begin{subfigure}[b]{0.48\linewidth}
    \centering
    \includegraphics[width=0.8\linewidth]{diodes_setup.png}
    \caption{Diodes: square (left) and rake (right)}
    \label{fig:diodes-setup}
  \end{subfigure}
  \caption{Measurement setup.}
  \label{fig:setup}
\end{figure}

\footnotetext{AI has been used to help create this figure}



Both devices have a contact on the top of the wafer, while the other contact is on the bottom of the wafer and is connected to the SMU. The square diode has a surface of $\SI{40\,000}{\micro \metre ^2}$ with a perimeter of $\SI{800}{\micro \metre}$ and the rake one $\SI{40\,320}{\micro \metre ^2}$ with a perimeter of $\SI{4072}{\micro \metre}$.

The wafer defines a minimal measurable current. We can measure it by extracting the current value on the log plot which is closest to zero. 
In theory the current should reach exactly zero, which is not the case in our measurement because of the limited resolution.

Extracting this value from the log plot (\cref{fig:iv-small-dark}) gives us:

\[
  I_{\mathrm{min}} \approx \SI{2,0e-8}{A}
\]


%=============================================================================
\subsection{I--V characteristics and operating regimes}\label{sec:iv-regimes}
In this part we will look at the I--V diagrams in more detail. At first, we display the diagrams in linear and log space. The log space is useful because the exponential
fall can be seen as linear. It also simplifies the identification of the minimal measurable current. From the measurements we get:
\begin{figure}[h!]
  \centering
    \centering
    \includegraphics[width=\textwidth]{../analysis/diode/data/plots.pdf}
    \caption{I--V characteristics of the diodes in linear (left) and logarithmic (right) scale}
    \label{fig:iv-small-dark}
\end{figure}

To explain the regions we first simplify the plot into a schematic. In this schematic we can identify five regions:

\begin{itemize}
  \item \textbf{series resistance:} at high forward bias the barrier $qV_{bi}$ of \cref{fig:diode-bands}
  is almost gone. The junction does not limit the current. The I--V relation becomes ohmic, i.e.\ linear with the slope $1/R_S$.
  \item \textbf{diffusion:} the forward bias lowers the barrier to $q(V_{bi}-V)$. Majority carriers can
  diffuse over it. We see an exponential rise.
  \item \textbf{generation-recombination:} at small forward bias only a few carriers can cross the barrier.
  Electrons and holes recombine in states due to the impurity of the crystal. This current rises
  slowly.
  \item \textbf{reverse leakage:} the reverse bias raises the barrier to $q(V_{bi}+V_R)$. 
  Only thermal minority carriers are carried across by the field. This current is the reverse current. This rises under illumination, because more carriers are created.
  \item \textbf{avalanche:} at large reverse bias the field in the depletion region is strong enough 
  for carriers to create new $e^-$--$h^+$ pairs by impact ionization. Therefore we get a self-sustaining carrier growth. The current rises strongly.
\end{itemize}

\begin{figure}[h!]
  \centering
  \begin{subfigure}[b]{0.62\linewidth}
    \centering
    \resizebox{\linewidth}{!}{
\begin{tikzpicture}[
  y=0.6cm,   % <-- hauteur réduite (1 = taille d'origine ; 0.5 = encore plus plat)
  >={Stealth[length=2.2mm]},
  declare function={
    Is  = 0.002;  % saturation current i_s (drawing units)
    nVT = 0.12;   % n k_B T / q
    Rs  = 0.45;   % series resistance (horizontal cm per vertical cm)
    Vb  = 4.0;    % junction voltage at which avalanche sets in
    Ib  = 0.002;  % avalanche current prefactor
    nb  = 0.12;   % steepness of the avalanche
    % current through the junction as a function of the junction voltage V_d
    id(\v) = Is*(exp(\v/nVT)-1) - Ib*exp(-(\v+Vb)/nb);
  },
  lbl/.style={font=\small},
  regime/.style={font=\small, text=black!55, align=center},
  sep/.style={densely dotted, black!35},
]

% region boundaries (V_Pos where the curve changes regime)
\draw[sep] (-0.40,-3.9) -- (-0.40,2.3);   % generation-recombination | diffusion
\draw[sep] (-0.95,-3.9) -- (-0.95,2.3);   % diffusion | series resistance
\draw[sep] ( 4.20,-3.9) -- ( 4.20,2.3);   % reverse leakage | avalanche

% axes
\draw[->] (-3.3,0) -- (6.6,0) node[right] {$V_{\mathrm{Pos}}$};
\draw[->] (0,-3.9) -- (0,2.5) node[above] {$I_{\mathrm{Pos}}$};


% real diode: V_a = V_d + R_s i_d  (Eq. 2), plotted as (-V_a, -i_d)
\draw[thick, blue!70!black, domain=-4.83:0.89, samples=600, variable=\v]
  plot ({-(\v + Rs*id(\v))}, {-id(\v)});

% regimes
\node[regime] at (-2.15,0.9) {series\\resistance};
\node[regime, rotate=90, anchor=west] at (-0.675,0.2) {diffusion};
\node[regime, rotate=90, anchor=west] at (-0.20,0.2) {G--R};
\node[regime] at (2.1,-0.6) {reverse leakage};
\node[regime] at (5.4,-0.6) {avalanche};

% forward / reverse bias
\draw[<->, black!55] (-3.3,-4.5) -- (-0.05,-4.5) node[midway, below, regime] {forward bias};
\draw[<->, black!55] (0.05,-4.5) -- (6.6,-4.5) node[midway, below, regime] {reverse bias};

\end{tikzpicture}
}
    \caption{Operating regimes}
    \label{fig:diode-regions}
  \end{subfigure}\hfill
  \begin{subfigure}[b]{0.34\linewidth}
    \centering
    % Band diagram of the n+/p junction in equilibrium.
% Energies in drawing units (1 = 1 cm). The p side is fixed, the n side is shifted by the
% applied voltage: up by qV in forward bias, down by qV_R in reverse bias.
\begin{tikzpicture}[
  >={Stealth[length=1.6mm]},
  font=\small,
  ec/.style={thick, blue!70!black},
  ev/.style={thick, red!70!black},
  ef/.style={dashed, black!70},
  lvl/.style={font=\footnotesize},
  carrier/.style={circle, inner sep=0pt, minimum size=4pt},
  el/.style={carrier, fill=black},
  ho/.style={carrier, draw=black, fill=white},
  ttl/.style={font=\small, align=center},
]

\def\Eg{1.2}    % band gap
\def\EcP{1.9}   % E_C on the p side (fixed)
\def\dEF{0.25}  % distance E_C - E_F (n side) and E_F - E_V (p side)

% #1 x offset, #2 shift of the n side (+qV forward, -qV_R reverse, 0 equilibrium),
% #3 depletion width, #4 barrier label, #5 panel title, #6 height of the barrier label
\def\W{4.6}    % panel width
\def\bandpanel#1#2#3#4#5#6{%
  \begin{scope}[xshift=#1cm]
    \pgfmathsetmacro\EcN{1.2+#2}
    \pgfmathsetmacro\EvN{\EcN-\Eg}
    \pgfmathsetmacro\EvP{\EcP-\Eg}
    \pgfmathsetmacro\EFn{\EcN-\dEF}
    \pgfmathsetmacro\EFp{\EvP+\dEF}
    \pgfmathsetmacro\xa{1.2}
    \pgfmathsetmacro\xb{1.2+#3}
    \pgfmathsetmacro\xm{(\xa+\xb)/2}
    % regions: n, depletion, p
    \fill[blue!15]   (0,-1.0)   rectangle (\xa,2.5);
    \fill[yellow!40] (\xa,-1.0) rectangle (\xb,2.5);
    \fill[red!15]    (\xb,-1.0) rectangle (\W,2.5);
    \node[font=\footnotesize] at (\xa/2,2.3) {n$^+$};
    \node[font=\footnotesize] at ({(\xb+\W)/2},2.3) {p};
    % conduction and valence band
    \draw[ec] (0,\EcN) -- (\xa,\EcN) .. controls (\xm,\EcN) and (\xm,\EcP) .. (\xb,\EcP) -- (\W,\EcP);
    \draw[ev] (0,\EvN) -- (\xa,\EvN) .. controls (\xm,\EvN) and (\xm,\EvP) .. (\xb,\EvP) -- (\W,\EvP);
    % (quasi-)Fermi levels
    \draw[ef] (0,\EFn) -- (\xb,\EFn);
    \draw[ef] (\xa,\EFp) -- (\W,\EFp);
    % barrier for electrons
    \draw[dotted] (\xa,\EcN) -- (\xb+0.35,\EcN);
    \draw[<->] (\xb+0.25,\EcN) -- (\xb+0.25,\EcP);
    \node[lvl, anchor=west] at (\xb+0.3,#6) {#4};
    \node[ttl, anchor=north] at (\W/2,-1.1) {#5};
  \end{scope}%
}

% ---------------------------------------------------------------- equilibrium
\bandpanel{0}{0}{1.0}{$qV_{bi}$}{}{1.55}
\node[lvl, left] at (0,1.2) {$E_C$};
\node[lvl, above] at (3.6,0.95) {$E_F$};
\node[lvl, left] at (0,0) {$E_V$};

\end{tikzpicture}

    \caption{Band diagram in equilibrium}
    \label{fig:diode-bands}
  \end{subfigure}
  \caption{(a) Operating regimes of the diode. (b) Band diagram of the n$^+$/p junction in
  equilibrium\protect\footnotemark}
  \label{fig:diode-regions-bands}
\end{figure}

\footnotetext{AI has been used to help create these figures}

\placeholder{colnclusion light / no light}
\pagebreak
%=============================================================================
\subsection{Extracted parameters: both modules, with and without light}\label{sec:iv-params}

In this section we want to measure characteristic values for the p-n diode. These consist of:

\begin{itemize}
  \item $\mathbf{V_\mathrm{\bf{th}}}$ \textbf{threshold voltage:} with rising forward bias, the junction barrier lowers. Below $V_\mathrm{th}$ the diode is fully forward biased and the slope is defined by the series resistance, above $V_\mathrm{th}$ the barrier still blocks most carriers. Therefore $V_\mathrm{th}$ marks the switch of the diode.
  We measure it by extending the left linear part of the curve to zero.
  \item $\mathbf{V_\mathrm{\bf{a}}}$ \textbf{avalanche voltage:} this is the voltage at which carriers gain enough kinetic energy to create new $e^- - h^+$ pairs. They will start a chain reaction. This is the point in the schematic where the current starts rising quickly. We extend the linearity of the rise and extract the root.

  \item $\mathbf{R_S}$ \textbf{series resistance:} this is the resistance that describes the ohmic losses outside of the junction. It is the resistance of the quasi-neutral $p$- and $n$-type regions. We measure it from the slope outside of the junction region in the I--V diagram.

  \item $\mathbf{I_R}$ \textbf{reverse current:} what we observe is that even in the reverse bias carriers flow through the junction. These are thermally generated minority carriers. We can measure the reverse current by measuring the current in the junction region of the diagram. 
\end{itemize}
From the measurement with no voltage we obtain that our cable has a leakage of $\SI{4}{\nano \A}$ which is negligible considering the current amplitudes used in the measurements.
\begin{figure}[h!]
  \centering
  
\begin{tikzpicture}[
  y=0.6cm,   % <-- hauteur réduite (1 = taille d'origine ; 0.5 = encore plus plat)
  >={Stealth[length=2.2mm]},
  declare function={
    Is  = 0.002;  % saturation current i_s (drawing units)
    nVT = 0.12;   % n k_B T / q
    Rs  = 0.45;   % series resistance (horizontal cm per vertical cm)
    Vb  = 4.0;    % junction voltage at which avalanche sets in
    Ib  = 0.002;  % avalanche current prefactor
    nb  = 0.12;   % steepness of the avalanche
    % current through the junction as a function of the junction voltage V_d
    id(\v) = Is*(exp(\v/nVT)-1) - Ib*exp(-(\v+Vb)/nb);
  },
  lbl/.style={font=\small},
  regime/.style={font=\small, text=black!55, align=center},
]

% axes
\draw[->] (-3.3,0) -- (6.6,0) node[right] {$V_{\mathrm{Pos}}$};
\draw[->] (0,-3.9) -- (0,2.5) node[above] {$I_{\mathrm{Pos}}$};


% real diode: V_a = V_d + R_s i_d  (Eq. 2), plotted as (-V_a, -i_d)
\draw[thick, blue!70!black, domain=-4.83:0.89, samples=600, variable=\v]
  plot ({-(\v + Rs*id(\v))}, {-id(\v)});

% forward: linear extrapolation of the high-current part -> V_th, slope 1/R_s
\draw[dashed, green!50!black] (-0.758,0) -- (-2.473,-3.5);
\fill[green!50!black] (-0.758,0) circle (1.4pt);
\node[lbl, green!50!black, above] at (-0.758,0.05) {$V_{\mathrm{th}}$};
\node[lbl, green!50!black, left] at (-1.85,-1.75) {slope $ s=1/R_s$};


% reverse: linear extrapolation of the avalanche branch -> V_av
\draw[dashed, red!70!black] (4.709,0) -- (5.882,2.3);
\fill[red!70!black] (4.709,0) circle (1.4pt);
\node[lbl, red!70!black, below] at (4.709,-0.05) {$V_{\mathrm{a}}$};
\node[lbl, red!70!black, right] at (5.45,1.35) {};

% regimes
\node[regime] at (6.0,-0.6) {avalanche};
\node[regime] at (2.4,-0.6) {reverse bias};
\node[regime] at (-2.1,0.55) {forward bias};

\end{tikzpicture}
  \caption[Extraction of values from diagram]{Extraction of values from diagram.\protect\footnotemark}
  \label{fig:diode}
\end{figure}

\footnotetext{AI has been used to help create this figure}
\begin{table}[H]
\centering\small
\caption{Parameters extracted from the I--V curves of the two diode modules.}
\label{tab:diode-iv}
\begin{tabular}{lcccccc}
\toprule
Module & Light & $V_{\mathrm{th}}$ (V) & $V_\mathrm{a}$ (V) & $R_\mathrm{S}$ (\si{\ohm}) & $I_\mathrm{R}(-1\,\mathrm V)$ ($\si{\micro \A}$) & $I_\mathrm{R}(-10\,\mathrm V)$ ($\si{\micro \A}$)\\
\midrule
Square & light & $-0.70$ & $18.70$ & $117.52$ & $47.60$ & $36.61$ \\
Square & dark & $-0.72$ & $18.79$ & $127.83$ & $3.01$  & $3.87$  \\
Rake   & light & $-0.66$ & $18.62$ & $126.65$ & $51.14$ & $55.46$ \\
Rake   & dark & $-0.70$ & $18.71$ & $126.74$ & $3.93$  & $4.31$ \\
\bottomrule
\end{tabular}
\end{table}

\placeholder{conclusion of the table}

%=============================================================================
\subsection{Ideality factor}\label{sec:iv-n}

In this section we want to calculate the ideality factor. This factor is the factor by which the voltage has to be increased
to achieve the behavior of the ideal diode. In the Boltzmann approximation the relation of the carrier population $n$ is given by:

\[
  n(E) \propto \exp\left(- \frac {E}{k_B T}\right) = \exp\left(- \frac {V_{\text{bi}} - V}{k_B T}\right) \propto \exp\left(\frac {V}{k_B T}\right)
\]

Because the ideality factor gives us the multiplier for $V$ we get:

\[
  I = I_0 \left(\exp\left(\frac {V}{ n k_B T}\right) - 1 \right)
\]

By now measuring the exponent $S$ by the slope of $I$ in the $\log$-scale, we can calculate $n$:

\begin{align*}
  S &= \frac{\Delta \ln(I)}{\Delta V}  &&= - \frac{\ln(\SI{1.004340e-5}) - \ln(\SI{8.526250e-7})}{(\SI{-0.5}{\V}) +\SI{0.1}{\V}}  \approx  \SI {6,17}{ \frac{1}{\V}} \\
  V_{\text{thermal}} &= \frac{k_B T }{q} &&= \frac{\SI{1.381e-23}{\joule\per\kelvin} \cdot \SI{300}{\kelvin}}{\SI{1.602e-19}{\coulomb}} \approx \SI{25,7}{ \milli \V} = \SI {0,0257} {\V} \\
  n &= \frac {1} {S \cdot V_\text{thermal}} &&= \frac{1}{\SI{6.17}{\per\V} \cdot \SI{0.0257}{\V}} \approx 6.32
\end{align*}

What we observe is a very unexpected value for $n$. Normally we would expect a value between $1$ and $2$.
Our value is far out of this range.

%=============================================================================
\subsection{Discussion}\label{sec:iv-discussion}
\placeholder{Item 7: dependence of the reverse current on light ($i_{ph}$) and on the
junction shape (area vs perimeter: rake vs square), values of $R_s$ and $n$, comparison
with an ideal diode.}

% ---------------------------------------------------------------------
\subsection{C--V characteristic of the junction (item 8)}\label{sec:diode-cv}
% ---------------------------------------------------------------------
\textbf{Setup.} The C--V curve was measured at \SI{1}{MHz} on the rake
diode, with the same setup as the MIS capacitor, from 0 to \SI{+15}{V} on the n contact
(step \SI{0.1}{V}). As the p substrate is on the chuck, a positive voltage on the n side
means \emph{reverse} bias: almost no current flows and we measure the capacitance of the
space-charge region (SCR). We did not use forward bias, because the current would be too large
to measure a meaningful capacitance.

\textbf{Result.} The capacitance decreases from \SI{30.1}{pF} at \SI{0}{V} to \SI{7.6}{pF} at
\SI{15}{V} (\cref{fig:diode-cv}a): when the reverse bias increases, the SCR becomes wider
($W=\varepsilon_{Si}A/C$ goes from \SIrange{0.14}{0.55}{\micro m}, with
$A=\SI{4.0e-4}{cm^2}$). For an n$^+$/p junction, almost all the SCR is on the p side and
\begin{equation}
\frac{1}{C^2}=\frac{2\,(V_{bi}+V_R)}{q\,\varepsilon_{Si}\,N_A\,A^2},
\label{eq:diode-cv}
\end{equation}
so the slope of $1/C^2$ versus $V_R$ gives the substrate doping $N_A$ and the intersection
with the voltage axis gives the built-in voltage $V_{bi}$ (same idea as the $1/C^2$ method of
\cref{sec:mis-extraction}). The parameters shown by MDC on this curve are not relevant, since the
software uses an MIS model.

\begin{figure}[htbp]
\centering
\begin{subfigure}[b]{0.42\linewidth}\centering
  \includegraphics[width=\linewidth]{mdc_diode_CV.jpg}\caption{}
\end{subfigure}\hfill
\begin{subfigure}[b]{0.40\linewidth}\centering
  \includegraphics[width=\linewidth]{diode_invC2.pdf}\caption{}
\end{subfigure}
\caption{(a) C--V curve of the n$^+$/p junction at \SI{1}{MHz} in reverse bias (MDC).
(b) $1/C^2$ as a function of the reverse bias: raw data (black) and after subtraction of a
parasitic capacitance $C_p=\SI{1.5}{pF}$ (red), with linear fits between 2 and \SI{15}{V}.}
\label{fig:diode-cv}
\end{figure}

\textbf{Extraction.} The raw $1/C^2$ curve is not perfectly straight. A linear fit gives
$N_A\approx\SI{7e16}{cm^{-3}}$ but $V_{bi}\approx\SI{1.7}{V}$, which is not possible (it should be
below \SI{1.1}{V}). We think that a small parasitic capacitance $C_p$ (contact pad and cables) is
measured in parallel with the junction. Its effect is larger when $C$ becomes small, which bends the
curve. Removing $C_p\approx\SI{1.5}{pF}$ makes the curve linear and gives
$N_A\approx\SI{4.4e16}{cm^{-3}}$ and $V_{bi}\approx\SI{0.83}{V}$, a reasonable value
(about \SI{0.9}{V} expected for an n$^+$/p junction with this doping). So this correction is
necessary (bonus question).

The doping is about 3 times higher than the one obtained with the MIS capacitor
(\SI{1.3e16}{cm^{-3}}, \cref{sec:mis-extraction}). A possible reason is that the real junction area
is larger than the drawn area, because the junction also extends on the sides of the n region;
since $N_A\propto1/A^2$, a small error on $A$ has a large effect. Comparing the square and rake
diodes would allow to check this.
